\chapter{Protocol blocks}
\label{ch:blocks}

\chapterbrief
  {Model the three ISO-14443-4 frame types as a small class hierarchy.}
  {Inheritance, polymorphism, and the factory method.}
  {A three-class hierarchy with tests, and one place where type is decided.}

\section{The problem}

Once a card is activated you exchange \textbf{T=CL blocks}. Every frame starts
with a PCB (Protocol Control Byte) whose top bits say which of three kinds it is:

\begin{center}\small
\begin{tabular}{@{}llll@{}}
\toprule
\thead{PCB bits} & \thead{Type} & \thead{Carries} & \thead{Purpose} \\
\midrule
\code{0xxxxxxx}  & \textbf{I-block} (Information) & your APDU & the actual data \\
\code{10xxxxxx}  & \textbf{R-block} (Receive ready) & nothing & acknowledgement \\
\code{11xxxxxx}  & \textbf{S-block} (Supervisory) & control info & ``wait'', ``deselect'' \\
\bottomrule
\end{tabular}
\end{center}

The C approach is a switch statement, and it works fine --- \emph{once}. The
problem is what happens over time. You need \code{describe()} for logging: one
switch. Then \code{expectsResponse()}: a second. Then \code{encode()}: a third.
Now three switches must be kept in agreement, and when a fourth block type
appears you must find all of them.

That is the smell: \textbf{the type of a thing is being re-decided at every place
you use it.}

\section{The idea: decide once, then let the object know what it is}

Make each block type a class. Decide which one you are holding \emph{exactly
once}, at the moment the byte arrives. Afterwards you hold an object that already
knows what it is, and you just call methods on it.

\begin{analogy}{The parcel depot}
A parcel arrives. You could sort it at every stage --- check ``is this fragile?''
at the truck, again at the belt, again at the van. Or you sort \emph{once} at
intake, put it in the right bin, and afterwards every stage handles ``whatever is
in this bin'' by that bin's rules.

One decision point, made once, near the door. The factory method is the intake
desk; polymorphism is the bin's rules applying automatically.
\end{analogy}

\begin{center}
\begin{tikzpicture}[font=\footnotesize\sffamily,
  b/.style={draw=brewernavy, rounded corners=2pt, inner sep=4pt,
            minimum width=30mm, minimum height=7.5mm},
  sub/.style={draw=brewergrey, rounded corners=2pt, inner sep=4pt,
            minimum width=27mm, minimum height=7.5mm}]
  \node[b] (base) {\code{ProtocolBlock}};
  \node[font=\scriptsize\itshape, color=brewergrey, right=2mm of base]
       {abstract; holds \code{pcb}};
  \node[sub, below left=13mm and 20mm of base] (i) {\code{InformationBlock}};
  \node[sub, below=13mm of base] (r) {\code{ReceiveReadyBlock}};
  \node[sub, below right=13mm and 20mm of base] (sb) {\code{SupervisoryBlock}};
  \foreach \n in {i, r, sb}{
    \draw[-{Triangle[open, length=5pt]}, brewernavy] (\n) -- (base);}
  \node[font=\scriptsize\itshape, color=brewergrey, below=1mm of i] {carries data};
  \node[font=\scriptsize\itshape, color=brewergrey, below=1mm of r] {ACK / NAK};
  \node[font=\scriptsize\itshape, color=brewergrey, below=1mm of sb] {control};
\end{tikzpicture}
\end{center}

\begin{concept}{Polymorphism}
\index{polymorphism}
Literally ``many shapes''. You call \code{block.describe()}. Java looks at what
the object \emph{actually is} at runtime --- I, R or S --- and runs that class's
version.

One call site, three behaviours, and no \code{if} anywhere.
\end{concept}

\section{Build it}

\subsection{The base}

\begin{jcode}
package com.midnightbrewer.hardware;

/**
 * One ISO-14443-4 (T=CL) block.
 *
 * All three types share the PCB, so that state lives here in the base -- which
 * is exactly why this is an abstract CLASS and not an interface (an interface
 * cannot hold the field).
 *
 * What DIFFERS is behaviour, so describe() is abstract: each subclass answers
 * it its own way.
 */
public abstract class ProtocolBlock {

    private final int pcb;

    // 'protected' -- only subclasses construct blocks, via the factory below.
    protected ProtocolBlock(int pcb) {
        this.pcb = pcb & 0xFF;
    }

    /** The raw PCB byte (0..255). */
    public int pcb() {
        return pcb;
    }

    /** Declared here, no body -- every subclass MUST answer. */
    public abstract String describe();

    // =================================================================
    // THE FACTORY -- the ONE place the type is decided.
    //
    // This switch exists once, here. Everywhere else you hold a
    // ProtocolBlock and call methods on it. No type checks, no instanceof.
    //
    //   b8 = 0         -> I-block
    //   b8 = 1, b7 = 0 -> R-block
    //   b8 = 1, b7 = 1 -> S-block
    // =================================================================
    public static ProtocolBlock from(int pcb) {
        // TODO 1: if bit 7 (0x80) is clear      -> new InformationBlock(pcb)
        // TODO 2: else if bit 6 (0x40) is clear -> new ReceiveReadyBlock(pcb)
        // TODO 3: else                          -> new SupervisoryBlock(pcb)
        return null;
    }
}
\end{jcode}

\begin{concept}{Why a static factory instead of \texttt{new}}
Because the caller does not know which class it wants --- that is precisely what
it is asking. \code{ProtocolBlock.from(0x02)} says ``give me whatever block this
byte represents''. The caller states intent; the factory decides the type.

This is the \textbf{Open/Closed Principle} in miniature: to add a block type you
add a class and one line in the factory. Every existing call site keeps working,
untouched. Compare the switch-everywhere version, where you must hunt through the
codebase and hope you found them all.
\end{concept}

\subsection{The three subclasses}

\begin{jcode}
package com.midnightbrewer.hardware;

public class InformationBlock extends ProtocolBlock {

    public InformationBlock(int pcb) {
        super(pcb);          // hand the PCB up to the base to store
    }

    @Override
    public String describe() {
        // TODO: "I-block (data), block number N"
        //       The block number is the low bit: pcb() & 0x01
        return null;
    }
}
\end{jcode}

\begin{jcode}
public class ReceiveReadyBlock extends ProtocolBlock {

    public ReceiveReadyBlock(int pcb) { super(pcb); }

    @Override
    public String describe() {
        // TODO: "R-block (ACK)" or "R-block (NAK)" -- bit 4 (0x10) set = NAK.
        return null;
    }
}
\end{jcode}

\begin{jcode}
public class SupervisoryBlock extends ProtocolBlock {

    public SupervisoryBlock(int pcb) { super(pcb); }

    @Override
    public String describe() {
        // TODO: bits 5..4 (mask 0x30) say which:
        //         0x00 -> "S-block (DESELECT)"
        //         0x30 -> "S-block (WTX -- card asking for more time)"
        //         else -> "S-block (unknown)"
        return null;
    }
}
\end{jcode}

\code{super(pcb)} calls the parent's constructor, and it must be the first
statement: the base must be fully built before the subclass adds to it.
Foundations before walls.

\begin{concept}{Always write \texttt{@Override}}
It is optional, and you should write it anyway. It tells the compiler ``I mean
to be replacing a parent method''. If you typo \code{describ()}, you get a
compile error instead of a method that silently never gets called. Free
insurance against a genuinely baffling class of bug.
\end{concept}

\section{Verify}

\begin{jcode}
package com.midnightbrewer.hardware;

import org.junit.jupiter.api.Test;
import static org.junit.jupiter.api.Assertions.assertEquals;
import static org.junit.jupiter.api.Assertions.assertTrue;

class ProtocolBlockTest {

    @Test
    void factoryPicksTheRightType() {
        assertTrue(ProtocolBlock.from(0x02) instanceof InformationBlock);
        assertTrue(ProtocolBlock.from(0x03) instanceof InformationBlock);
        assertTrue(ProtocolBlock.from(0xA2) instanceof ReceiveReadyBlock);
        assertTrue(ProtocolBlock.from(0xF2) instanceof SupervisoryBlock);
    }

    @Test
    void polymorphicDescribe() {
        // Same variable type, same method call, three different answers.
        ProtocolBlock[] blocks = {
            ProtocolBlock.from(0x02),
            ProtocolBlock.from(0xA2),
            ProtocolBlock.from(0xF2),
        };
        assertTrue(blocks[0].describe().startsWith("I-block"));
        assertTrue(blocks[1].describe().startsWith("R-block"));
        assertTrue(blocks[2].describe().startsWith("S-block"));

        assertEquals(0, ProtocolBlock.from(0x02).pcb() & 0x01);
        assertEquals(1, ProtocolBlock.from(0x03).pcb() & 0x01);
    }
}
\end{jcode}

Look closely at \code{polymorphicDescribe}: the array's type is
\code{ProtocolBlock[]}. The code calling \code{describe()} has \textbf{no idea}
which subclass each element is, and does not need to. That is the entire point.

\begin{deeper}{\texttt{instanceof} in tests versus in production}
Using \code{instanceof} in the first test is fine --- a test's job is to verify
the factory picks correctly.

\code{instanceof} in \emph{production} code is usually a signal that you should
have added a method to the base instead. If you find yourself writing
\code{if (b instanceof InformationBlock) \{ ... \}}, ask what question you are
really asking, and give the base class a method that answers it.
\end{deeper}

\begin{verifybox}
\begin{shell}
mvn test -Dtest=ProtocolBlockTest
\end{shell}
\end{verifybox}

\section{Inheritance versus composition}
\index{composition}\index{inheritance}

You have now used both ways one class can relate to another:

\begin{itemize}
  \item \textbf{Composition (HAS-A)} --- \code{Rc522Reader} \emph{has a}
        \code{Rc522Driver}; \code{Rc522Driver} \emph{has a} \code{SpiLink}.
  \item \textbf{Inheritance (IS-A)} --- \code{Rc522Reader} \emph{is a}
        \code{Iso14443Reader}; \code{InformationBlock} \emph{is a}
        \code{ProtocolBlock}.
\end{itemize}

The test is whether ``IS-A'' is literally true. An I-block \emph{is a} protocol
block --- yes. Is a reader \emph{a kind of} driver? No: a reader \emph{uses} a
driver. Different relationships entirely.

\begin{trap}{The \texttt{extends Pi4jSpiLink} disaster}
Writing \code{class Rc522Driver extends Pi4jSpiLink} compiles fine and looks
like a shortcut --- now the driver can call \code{transfer} directly!

Then your unit test constructs a driver, which constructs a Pi4J context, which
hunts for a GPIO chip that does not exist on your laptop, and the test dies with
\code{No gpio chip found}. You have welded the driver to the hardware and
destroyed the seam you built in Chapter~\ref{ch:spi}.

\textbf{Default to composition.} Reach for inheritance only when you genuinely
have a family sharing an algorithm --- like these blocks, or the reader base.
\end{trap}

\section{Commit}

\begin{shell}
git add -A
git commit -m "M4: ProtocolBlock hierarchy -- I/R/S blocks with a factory"
\end{shell}
